% ch-a6-05.tex —— 第 6 篇（Feng–Harris）续：PDF 532–543 / 印刷 519–530
% 前一片 ch-a6-04.tex 覆盖印刷 507–518；本片从印刷 p.519（§7.2 末尾）起，
% 经 §7.3、§7.4、§8.1–§8.4，止于 p.530 页末（§8.4.1 开头）。
% 衔接：p.531 末句为 "…where C runs through possible fields"，故 a6-04 末句已把 "of coefficients" 一并译出，本片不再重译该残句；
% 若上一片已把 "of coefficients" 一并译出，请删去下页首的「……的系数。」五字。

\ysrc{519}关键点在于\footnote{通常把 $\varpi_\pi$ 定义为那些能延拓到某个适当整模型上的类所构成的空间。在应用中，Venkatesh 解释了这一条件为何可以忽略。例如可参看 [PV21, p. 4]。}

\begin{conj}{7.2.6}（Beilinson, Bloch--Kato）\label{thm:a6-05-1}
\begin{enumerate}[(i)]
\item 空间 $\varpi_\pi$ 的维数为 $\ell_0$。
\item Beilinson 正则映射给出同构
\[
\varpi_\pi\otimes_\Q\C\isom H^1_D(M_{\pi,\mathrm{HdR}},1),
\]
其中 $H^1_D(*,i)$ 是 Deligne 上同调，$M_{\pi,\mathrm{HdR}}$ 表示 Hodge--de Rham 实现。
\item 对 $M_\pi$ 具有好约化的任一素 $p$，$p$-进正则映射给出同构
\[
\varpi_\pi\otimes_\Q\Q_p\isom H^1_f(\Q,M_\pi(1)),
\]
其中 $H^1_f$ 是 Bloch--Kato Selmer 群。
\end{enumerate}
\end{conj}

每个猜想还断言：上述同构右端的空间维数为 $\ell_0$。这往往可以由其他猜想证明，或从其他猜想推出。对我们的目的而言最重要的是：这些猜想在 $\wedge^\bullet(\mathfrak a^*)$ 上给出有理结构，而借助 (7.2.7) 可以把该结构与上同调空间 $H^*_0(S(G,X),\widetilde V_\C)[\pi_f]$ 上由拓扑上同调诱导的有理结构作比较。具体地说，我们有

\begin{fact}{7.2.7}（[PV21], (5.1.1)）\label{thm:a6-05-2}
猜想 7.2.6 的 (i) 与 (ii) 给出典范同构
\[
\varpi_\pi\otimes_\Q\C\isom\mathfrak a^*.
\]
\end{fact}

于是猜想 7.2.6 赋予 $\mathfrak a^*$ 一个典范的 $\Q$-有理结构。

[PV21, Ven19, GV18] 中的诸猜想都回到这个 $\Q$-有理结构，尤其回到下述猜想。由于 $\pi_f$ 在上同调中得到实现，它在某个数域 $E(\pi_f)$ 上有模型。\footnote{若 $G=\mathrm{GL}(n)$，则由重数一定理可知，该域可以认定为 $\pi_f$ 的有理性域，即 $\Gal(\overline{\Q}/\Q)$ 中固定 $\pi_f$ 的同构类的子群所对应的固定域。一般地，$E(\pi_f)$ 可能略大一些。}于是 $H^*_0(S(G,X),\widetilde V_\C)$ 的 $\pi_f$-同型分量是 $E(\pi_f)$-有理上同调 $H^*_0(S(G,X),\widetilde V_{E(\pi_f)})$ 的直和项。下面陈述 Prasanna 与 Venkatesh 主猜想的一个轻微推广：

\begin{conj}{7.2.8}（[PV21], Conjecture 1.2.1）\label{thm:a6-05-3}
$\wedge^\bullet\varpi_\pi$ 在 $H^*_0(S(G,X),\widetilde V_\C)[\pi_f]$ 上的作用保持 $E(\pi_f)$-有理结构。
\end{conj}

关于 $\varpi_\pi$ 的 $p$-进实现的作用，我们在 §8 中更详细地描述。

\subsection{志村簇的凝聚上同调}

我们回到 §7.2 的记号，但现在假设 $X$ 带有 $G(\R)$-不变的 Hermite 结构。我们始终假设 $Z_G$ 平凡；这当然没有必要，但可简化叙述。定义
\[
\mathfrak P=\mathfrak P_h=\mathfrak k\oplus\mathfrak p^-,
\]
记号同 (7.2.2)。这是以 $\mathfrak p^-$ 为幂幺根的极大抛物子代数。于是局部对称空间 $S(G,X)$ 是志村簇，并且

\ysrc{520}它在数域 $E(G,X)$ 上有典范模型，该模型在 $G(\mathbb A_f)$ 的右作用下保持。这就诱导出各个有限水平商 $S_{K_f}(G,X)$ 上的典范模型。

设 $\tau:K_\infty\to\mathrm{GL}(W_\tau)$ 是不可约表示，定义全纯向量丛
\begin{equation}\tag{7.3.1}\label{eq:a6-05-1}
\mathcal E_\tau=\varprojlim_{K_f}G(\Q)\backslash G(\mathbb A)\times W_\tau/K_\infty\times K_f,
\end{equation}
其中现在 $G(\Q)$ 从左边作用在 $G(\mathbb A)$ 上，$K_f$ 从右边作用在 $G(\mathbb A_f)$ 上，$K_\infty$ 从右边作用在 $G(\R)$ 上、从左边作用在 $W_\tau$ 上。若 $K_f$ 充分小，则丛 $\mathcal E_\tau$ 下降为 $S_{K_f}(G,X)$ 上的向量丛（仍记作 $\mathcal E_\tau$），它在 $E(G,X)$ 的有限扩张 $E(G,X,\tau)$ 上有典范模型。实践中我们固定水平子群 $K_f$，使 $S_{K_f}(G,X)$ 具有典范整模型，如 Kisin、Shin 与 Zhu [KSZ21] 的工作所示；这预先假定 $S(G,X)$ 是阿贝尔型志村簇，而后者与带附加结构的阿贝尔簇的模空间密切相关。\footnote{校样附注：本文写就之后，Bakker、Shankar 与 Tsimerman 在特征充分大的情形下构造了例外志村簇的典范整模型。}

暂时假设 $G$ 在 $\Q$ 上各向异性，于是 $S(G,X)$ 是（pro-）投射的。以（附着于）向量丛 $\mathcal E_\tau$ 的凝聚层为系数的 $S(G,X)$ 上同调，可以用微分形式（Dolbeault 上同调）表达；后者又可以用相对 Lie 代数上同调写出，类比于 Matsushima 公式：
\begin{equation}\tag{7.3.2}\label{eq:a6-05-2}
H^*(S(G,X),\mathcal E_\tau)\isom H^*(\mathfrak P,\mathfrak k;A(G)\otimes W_\tau)\isom\bigoplus_\pi m(\pi)H^*(\mathfrak P,\mathfrak k;\pi_\infty\otimes W_\tau)\otimes\pi_f.
\end{equation}
（同样，应把 $(\mathfrak P,\mathfrak k)$-上同调换成 $(\mathfrak P,K_h)$-上同调，其中 $K_h\supset K_\infty$ 是 $h$ 在 $G(\R)$ 中的稳定化子，未必连通。）当 $S(G,X)$ 非投射时，需要把 (7.3.2) 左端换成 $H^*(S(G,X)_\Sigma,\mathcal E_\tau^{\mathrm{can}})$，其中 $S(G,X)_\Sigma$ 是选取得当的（光滑投射）环面紧化，$\mathcal E_\tau^{\mathrm{can}}$ 是 Mumford（或 Deligne）的典范延拓，把 $\mathcal E_\tau$ 延拓成 $S(G,X)_\Sigma$ 上的向量丛；结果与 $\Sigma$ 的选取无关。尖点上同调 (7.2.6) 的类比
\begin{equation}\tag{7.3.3}\label{eq:a6-05-3}
H^*_0(S(G,X),\mathcal E_\tau)\isom\bigoplus_{\pi\subset A_0(G)}m(\pi)H^*(\mathfrak P,\mathfrak k;\pi_\infty\otimes W_\tau)\otimes\pi_f
\end{equation}
曾在 [Har90] 中研究；在多数情形下，其左端用 $\mathcal E_\tau$ 到 $S(G,X)_\Sigma$ 的两个不同延拓按代数方式定义。Franke 定理 (7.2.4) 的类比对凝聚上同调也成立，最近由 Jun Su 的学位论文 [Su] 证明。

设 $\tau$ 作为 $\mathfrak k$ 的包络代数的表示，其无穷小特征标限制到 $\mathfrak g$ 的包络代数上时，等于某个有限维不可约表示的无穷小特征标。此时 Schmid 定理断言：存在唯一的 $\pi_\infty$ 与唯一的 $q(\tau)\in\{0,1,\dots,\dim X\}$，使得 $H^{q(\tau)}(\mathfrak P,\mathfrak k;\pi_\infty\otimes W_\tau)\ne0$；且该上同调空间的维数为 1，$\pi_\infty$ 是离散系列。这一理论在 [Har90] 中有详尽叙述。特别地，整体凝聚上同调 $H^q_0(S(G,X),\mathcal E_\tau)$ 仅在 $q=q(\tau)$ 时非平凡。

\ysrc{521}Venkatesh 的动机猜想针对凝聚上同调有一个版本，将在 §8.3 中讨论；它涉及使 $H^*_0(S(G,X),\mathcal E_\tau)$ 在多个次数上都有上同调的 $\tau$，相应的 $\pi_\infty$ 属于离散系列的非退化极限（NLDS）。[BHR94] 在 Floyd Williams 工作的基础上证明：任一给定的 NLDS $\pi_\infty$ 只对唯一的 $q$ 贡献到 $H^q(\mathfrak P,\mathfrak k;\pi_\infty\otimes W_\tau)$。设 $P(\tau)$ 表示使 $H^*(\mathfrak P,\mathfrak k;\sigma\otimes W_\tau)\ne0$ 的 NLDS $\sigma$ 的集合，并定义
\[
\Pi(\tau)=\oplus_{\sigma\in P(\tau)}\sigma.
\]
就这一定义而言，最方便的做法是把 $G(\R)$ 的任一 NLDS 限制到单位连通分支后得到的各个分量，当作彼此独立的表示 $\sigma$ 来处理。

\begin{ques}{7.3.1}\label{thm:a6-05-4}
是否存在 $\mathfrak p$ 的阿贝尔 Lie 子代数 $\mathfrak a$，使得 $H^*(\mathfrak P,\mathfrak k;\Pi(\tau)\otimes W_\tau)$ 是 $\wedge^\bullet(\mathfrak a^*)$ 上的自由模？
\end{ques}

在 $G=\mathrm{GL}_2$ 这一（公认不具代表性的）情形，答案是肯定的；但这一问题似乎尚未在更一般的情形中得到考察。

\subsection{$q$-进导出 Hecke 代数在上同调上的作用}

我们回到 §7.2 的拓扑设定，具体回到 Prasanna--Venkatesh 猜想 7.2.7。故设 $\pi_f$ 如该节所述；特别地，我们假设所有 $\pi_\infty\in[\pi_f]_\infty$ 都是缓增的。我们假设 $E(\pi_f)=\Q$；这一假设容易放松，代价是记号变复杂。我们同样始终假设表示 $V$ 定义在 $\Q$ 上，并记 $\widetilde V$ 为相应的以 $\Q$ 为系数的局部系统。

设 $p$ 是使局部分量 $\pi_p$ 非分歧的素数。我们固定水平子群 $K_0=\prod_v K_v\subset G(\mathbb A_f)$，其中 $K_p$ 是超特殊极大紧子群，使得 $\pi_f^{K_0}\ne0$。我们将使用局部对称空间 $S_{K_0}(G,X)$ 的 $\Q_p$-值上同调，系数在 $\widetilde V\otimes\Q_p$ 中。更确切地说，我们选取 $\Z_p$-局部系统 $\Lambda_V\subset\widetilde V$，并如 [Ven19, §6.1] 那样把 Hecke 代数 $\mathbb T_{K_0}$ 定义为由好素数处的 Hecke 算子生成的 $\Z_p$-代数，这些算子作用为链复形 $\Lambda_V$ 的自同态，而链复形看作导出范畴中的对象。于是 $\pi_f$ 决定一个极大理想 $\mathfrak m=\mathfrak m_\pi\subset\mathbb T_{K_0}$，我们留待后用。

记 $S_{K_0}=S_{K_0}(G,X)$。对任意 $r=1,2,\dots$，令 $\Lambda_{V,r}=\Lambda_V/p^r\Lambda_V$，$A_r=\Z/p^r\Z$。我们把对象
\[
R\Gamma(S_{K_0},\Lambda_{V,r})
\]
看作 $\mathbb T_{K_0}$-模。设 $q$ 是使 $K_q$ 为超特殊极大紧子群的素数。考虑由包含 $\prod_{v\ne q}K_v$ 的开紧子群 $U\subset K_0$ 构成的族 $\mathcal F_q$；对每个 $U\in\mathcal F_q$，考虑链复形 $C^\bullet(S_U(G,X),\Lambda_{V,r})$，视作 $A_r$-模范畴的导出范畴中的对象。于是 $G_q=G(\Q_q)$ 作用在 $C^\bullet_q(A_r)=\varinjlim C^\bullet(S_U(G,X),\Lambda_{V,r})$ 上。设 $\mathcal U_q$ 是泛 $A_r[G_q]$-模 $\mathcal U_{A_r}(K_q)$。则正如 §6.2 中那样，Frobenius 互反律给出等同
\begin{equation}\tag{7.4.1}\label{eq:a6-05-4}
R\Hom_{A_r[K_q]}(A_r,C^\bullet_q(A_r))\isom R\Hom_{A_r[G_q]}(\mathcal U_{A_r}(K_q),C^\bullet_q(A_r)),
\end{equation}
从而由事实 6.2.1，导出 Hecke 代数 $\mathcal H^*_{A_r}(G_q,K_q)$ 作用在 $H^*(K_q,C^\bullet_q(A_r))$ 上。

\ysrc{522}由 [HV19, §7.5] 的讨论，可以把超上同调 $H^*(K_q,C^\bullet_q(A_r))$ 与 $H^*(S_{K_0},\Lambda_{V,r})$ 等同起来。于是形式上我们有

\begin{prop}{7.4.1}\label{thm:a6-05-5}
对任意 $r$ 与任意 $q$，导出 Hecke 代数 $\mathcal H(G_q,K_q;A_r)$ 典范地作用在 $H^*(S_{K_0},\Lambda_{V,r})$ 上。
\end{prop}

命题 7.4.1 中的作用只在 $\mathcal H_{A_r}(G_q,K_q)$ 非平凡时才有意义。故设 $q\equiv1\pmod{p^r}$。对这种 $q$，由定理 6.1.6 已经看到
\begin{equation}\tag{7.4.2}\label{eq:a6-05-5}
\mathcal H^*(G_q,K_q;A_r)\isom\mathcal H^*(T(\Q_q),T(\Z_q);A_r)^W\isom(A_r[X^*(T)]\otimes_{A_r}H^*(T_q;A_r))^W,
\end{equation}
其中 $W$ 表示 $G$ 的 Weyl 群，右端代数的记号作了相应调整。

\begin{rem}{7.4.2}\label{thm:a6-05-6}
当 $q\equiv1\pmod{p^r}$ 时，$\mathcal H^1_{A_r}(G_q,K_q)$ 中单个导出 Hecke 算子在模 $p^r$ 上同调上的作用可以显式写出：它用从 Iwahori 水平过渡到 $q$ 处的 pro-$q$ Iwahori 水平时、次数为 $p^r$ 的循环非分歧拓扑覆盖的模 $p^r$ 特征类来表达。于是这些算子可看作钻石算子的导出版本，而钻石算子在 Taylor--Wiles 方法中起核心作用。在模曲线的凝聚上同调情形，§8.3 对此有更详细的解释。拓扑版本由 [Ven19] 发展，二者完全类似。
\end{rem}

Venkatesh 在 [Ven19, §2.13] 中解释了如何取上述 $q$ 的集合的递增并，同时在 $r$ 上取极限，从而定义作用于 $H^*(S_{K_0},\Lambda_V)$ 的整体 $p$-进导出 Hecke 代数的各种版本。我们选取严格整体导出 Hecke 代数，即 [Ven19] 中所记的 $\widetilde{\mathbb T}'$，并假设 $p$ 与 Weyl 群的阶互素。于是 $\widetilde{\mathbb T}'$ 是分次交换的，且包含除 $p$ 以外的好素数处的非导出 Hecke 算子。这些算子生成通常的（非导出）Hecke 代数 $\mathbb T_{K_0}\subset\widetilde{\mathbb T}'$。此外，$\widetilde{\mathbb T}'$ 的作用保持 $\mathbb T_{K_0}$ 的极大理想处的局部化 $H^*(S_{K_0},\Lambda_V)_{\mathfrak m}=H^*(S_{K_0},\Lambda_V)_{\mathfrak m_\pi}$，这个极大理想我们一直留到此刻才用。

[Ven19] 的第一个主要结果是

\begin{thm}{7.4.3}（[Ven19, Theorem 7.6]）\label{thm:a6-05-7}
在下面列出的假设下，$H^*(S_{K_0},\Lambda_V)_{\mathfrak m}$ 作为 $\widetilde{\mathbb T}'$ 上的分次模，由它的最低次分量 $H^{q_0(G)}(S_{K_0},\Lambda_V)_{\mathfrak m}$ 生成。
\end{thm}

其证明把 Calegari--Geraghty 的想法（Khare 与 Thorne 发展的版本）改写成导出结构的语言。证明依赖 [Ven19, §6.1] 中列出的技术性假设，以及 [Ven19, §6.2] 中的一个实质性假设。最重要的技术性假设是：$\mathfrak m$-局部化同调 [N.B.] $H_*(S_{K_0},\Lambda_V)_{\mathfrak m}$ 集中在区间 $[q_0(G),q_0(G)+\ell_0(G)]$ 内——这个区间我们已在缓增表示的相对 Lie 代数上同调计算中见过。与此相关的是 $H_*(S_{K_0},\Lambda_V)$ 无挠这一假设。最后，假设 $G$ 是分裂 $\Q$-群；这似乎排除了对全实域上的单群的推广，例如，尽管类似的定理对这类群也必定成立。

更严重的假设是 $\pi$ 配有 $p$-进伽罗瓦参数。这是一个连续同态
\begin{equation}\tag{7.4.3}\label{eq:a6-05-6}
\rho_\pi:\Gal(\overline{\Q}/\Q)\to\check G(\overline{\Q}_p),
\end{equation}
\ysrc{523}其中 $\check G$ 是 $G$ 的 Langlands 对偶群，它在非分歧素 $q$ 处的 Frobenius 共轭类对应于 $\pi_q$ 的 Satake 参数，而它到 $p$ 处分解群上的限制在适当意义下是晶体的。这样的 $\rho_\pi$ 的存在性，当 $G$ 为实域或 CM 域上的一般线性群时是已知的——前提是我们关于 $\pi$ 贡献到某个 $S_{K_0}$ 的上同调的假设。Venkatesh 接着对约化 $\overline{\rho}_\pi\bmod p$ 的像作假设，这些假设在 Taylor--Wiles 理论中很常见；特别是，$\overline{\rho}_\pi$ 的像不包含在 $\check G$ 的真抛物子群中，且不太小。

\begin{rem}{7.4.4}\label{thm:a6-05-8}
严格说来，Venkatesh 只在 $V$ 是平凡表示时证明定理 7.4.3。此外，他假设 $\pi$ 在水平 $K_0$ 处不同余于任何其他表示，并且 $\mathbb T_{K_0,\mathfrak m}\isom\Z_p$。据我们所知，还没有人在不设这些条件的情况下写出证明；这些条件大概需要对 $p$ 加额外限制。
\end{rem}

第二个主要结果是构造动机上同调的 $p$-进实现的作用（如猜想 7.2.6 (iii)），它作用在 $H^*(S_{K_0},\Lambda_V)_{\mathfrak m}$ 上，实现为把次数增加 1 的自同态。这将在下一节讨论；下一节专门讨论 Venkatesh 关于这一作用的主猜想，以及它在其他设定下的推广。

\section{Venkatesh 的动机猜想}

\subsection{动机上同调的作用}

我们保持 §7.4 的假设。于是 Venkatesh 的定理 7.4.3 告诉我们：局部对称空间的 $\ell$-进上同调，在对应于 $\pi$ 的极大理想处局部化后，是严格整体导出 Hecke 代数（DHA）上的循环模。另一方面，解析理论告诉我们：与 $\Q$ 张量之后，它是维数为 $\ell_0(G)$ 的向量空间的外代数上的循环模。这一不变量在严格整体 DHA 的定义中不起作用，但通过 Bloch--Kato 猜想与动机上同调相联系。$\ell_0(G)$ 的这两次出现之间的关系，由下述典范映射给出：从猜想 7.2.6 (iii) 右端的对偶伴随 Selmer 群，到严格整体 DHA 的次数 1 部分，该映射由一条局部互反律刻画。

不过，这条局部互反律是在整体语境中陈述的。我们把 $\pi$ 与 (7.4.3) 中的伽罗瓦表示实现在 $p$-进域的整数环 $\OO$ 上，剩余域为 $k$。为 $\pi$ 选取一个 Taylor--Wiles 系统，即一组满足 $q_i\equiv1\pmod{p^r}$ 的素数 $Q_r=(q_1,\dots,q_m)$，如 [Ven19, §6.3] 所述；相关性质参见该处。\footnote{另见本文集中 Caraiani--Shin 与 Emerton--Gee--Hellmann 的各章。实践中当然可以在 $\Q$ 以外的基域上工作，此时 $q_i$ 将是素理想。}记 $\mathcal H^*_{q,A_r}$ 为 $\mathcal H^*(G_q,K_q;A_r)$。特别地，$\pi$ 在每个 $q\in Q_r$ 处非分歧，且每个 $\pi_q$（$q\in Q_r$）的 Satake 参数强正则。这意味着：若 $\mathfrak m_q\subset H(T(\Q_q),T(\Z_q);A_r)^W$ 是由极大理想 $\mathfrak m_\pi$ 诱导的极大理想，则在
\[
\Spec\bigl(H(T(\Q_q),T(\Z_q),A_r)=A_r[X^*(T^\vee)]=A_r[X_*(T)]\bigr)
\]
中，$\mathfrak m_q$ 之上的纤维由 $|W|$ 个极大理想构成。我们在
\[
\Spec\bigl(\mathcal H^0(T(\Q_q),T(\Z_q);A_r)\bigr)
\]
中各选取一个极大理想，

\ysrc{524}对每个 $q$ 都如此取；这对应选取 Langlands 对偶环面 $T^\vee(k)$ 中一个落在 $\overline{\rho}_\pi(\mathrm{Frob}_q)$ 的共轭类内的元素，其中 $\overline{\rho}$ 表示 $\rho_\pi$ 在 $G^\vee(k)$ 中的约化。由于 $W$ 作用在 (7.4.2) 中张量积的两个因子上，这一选取连同导出 Satake 同构使我们能定义典范单射 [Ven19, §8.17]
\begin{equation}\tag{8.1.1}\label{eq:a6-05-7}
\iota_{q,r}:H^1(T_q;A_r)\hookrightarrow\mathcal H^1_{q,A_r}.
\end{equation}

我们可以在对偶伴随 Selmer 群 $H^1_f(\Q,M_\pi(1))$ 中定义一个 $\OO$-格。想到猜想 7.2.6，我们令 $\varpi^*_{\pi,p}$ 表示这个格的 $\OO$-对偶。在目前的假设下，Venkatesh 在 [Ven19, Lemma 8.9] 中证明 $\varpi_{\pi,p}$ 是 $\OO$ 上的自由模，秩为 $\ell_0(G)$。注意，整体对偶把 $\varpi^*_{\pi,p}$ 与（阻碍空间）$H^2_f(\Q,\mathrm{Ad}(\rho_\pi))$ 中的一个格等同起来。

还有第二个映射
\begin{equation}\tag{8.1.2}\label{eq:a6-05-8}
f_{q,r}:H^1(T_q,A_r)\to\varpi^*_{\pi,p}/p^n,
\end{equation}
它由伽罗瓦上同调的局部对偶定义（见 [Ven19] 第 78 页）。

Venkatesh 的第二个定理（在定理 7.4.3 之后）是下述结果的一个更精确版本：

\begin{thm}{8.1.1}（[Ven19, Theorem 8.5]）\label{thm:a6-05-9}
存在 $\varpi^*_{\pi,p}$ 在 $H^*(S_{K_0},\Lambda_V)_{\mathfrak m}$ 上的作用，实现为次数 $+1$ 的自同态。该作用由下述性质唯一刻画：对任意 $r\ge1$ 与任意 Taylor--Wiles 素 $q\in Q_{a(r)}$（$a(r)$ 充分大），$H^1(T_q,A_r)$ 在 $H^*(S_{K_0},\Lambda_V)_{\mathfrak m}$ 上的两个作用——经由导出 Hecke 算子（通过 $\iota_{q,r}$）与经由 $\varpi^*_{\pi,p}$ 的作用（通过 $f_{q,r}$）——彼此吻合。

特别地，外代数 $\wedge^\bullet\varpi^*_{\pi,p}$ 嵌入到 $H^*(S_{K_0},\Lambda_V)_{\mathfrak m}$ 的自同态构成的严格整体 DHA $\widetilde{\mathbb T}'$ 中。
\end{thm}

\begin{rem}{8.1.2}（谱 Hecke 代数）\label{thm:a6-05-10}
定理 8.1.1 的证明是 [Ven19] 中一项密集的计算。另一种视角由 [Fen20] 给出。该文构造了一个叫作谱 Hecke 代数的对象；它是作用在 Langlands 对应的谱（即伽罗瓦）侧的导出 Hecke 代数。这一构造的动机来自几何 Langlands，我们在此不解释该动机；细节参见 [Fen20, Introduction]。
\end{rem}

$G$ 的通常球 Hecke 代数，是 $G(\Z_q)\backslash G(\Q_q)/G(\Z_q)$ 上（看作双陪集集合）的函数空间。自然地把这个双陪集空间看作群胚，即记住自同构。于是它具有非平凡的拓扑结构，而取上同调（而非只取函数）就得到导出 Hecke 代数。对 $G=\mathrm{GL}_n$，该群胚可以解释为三元组 $(F_1,F_2,\tau)$ 构成的空间，其中 $F_1,F_2$ 是 $\Spec\Z_q$ 上秩 $n$ 的向量丛，$\tau$ 是它们到 $\Spec\Q_q$ 上的限制之间的等同。

为构造谱 Hecke 代数，我们改为考虑三元组 $(F_1,F_2,\tau)$ 构成的空间，其中 $F_1,F_2$ 是 $\Spec\Z_q$ 上秩 $n$ 的局部系统，$\tau$ 是它们到 $\Spec\Q_q$ 上的限制之间的等同。两种情形的语言虽然平行，几何却大不相同。在第二种情形，只有按导出意义作构造才有意思，因为 $\Spec\Z_q$ 上秩 $n$ 局部系统的空间嵌入到 $\Spec\Q_q$ 上秩 $n$ 局部系统的空间之中（于是 $F_2$ 的信息是多余的）。我们称这个导出空间为「谱 Hecke 叠」，它的函数本质上就是 [Fen20] 的谱 Hecke 代数。

谱 Hecke 代数在导出伽罗瓦形变环上有一个余作用，[Fen20] 算出它与 Venkatesh 关于导出 Hecke 代数作用的互反律对偶。

\ysrc{525}此外，[Fen20] 说明谱 Hecke 代数上有一个余乘法，使它和导出 Hecke 代数对偶。然后 [Fen20] 给出 Venkatesh 互反律的一个概念性解释：它是在导出 Hecke 代数的作用与谱 Hecke 代数的作用之间的「导出的局部--整体相容性」。

现在我们把动机上同调与 $S_{K_0}$ 的上同调都与 $\Q_p$ 张量，并假设 Beilinson--Bloch--Kato 猜想 7.2.6 (iii)。令 $\varpi^*_{\pi,E(\pi_f)}\subset\varpi^*_{\pi,p}\otimes\Q_p$ 为那些与动机上同调的配对落在 $E(\pi_f)$ 中的类构成的空间。记 $H^*(S(G,X),\widetilde V)_\pi$ 为 $H^*(S(G,X),\widetilde V)$ 的 $\pi_f$-同型子空间。Venkatesh 动机猜想的 $p$-进版本是

\begin{conj}{8.1.3}（[Ven19, Conjecture 8.8]）\label{thm:a6-05-11}
在 $\wedge^\bullet\varpi^*_{\pi,p}$ 通过定理 8.1.1 定义的作用之下，
\[
[H^*(S(G,X),\widetilde V)_\pi]^{K_0}\otimes\Q_p\subset H^*(S_{K_0},\Lambda_V)_{\mathfrak m}\otimes\Q_p
\]
上 $\varpi^*_{\pi,E(\pi_f)}$ 的作用保持由 $H^*(S(G,X),\Lambda_V)_\pi$ 诱导的 $E(\pi_f)$-有理性结构。
\end{conj}

除 [DHRV22] 关于模曲线凝聚上同调的一个类似猜想的结果（见 §8.3）之外，支持这一猜想的最强证据是 [PV21] 为复上同调的类似猜想提供的大量证据。在那个猜想中，取值于 Bloch--Kato Selmer 群的 $p$-进正则映射换成取值于 Deligne 上同调的 Beilinson 正则映射，如猜想 7.2.6 (ii) 那样。由于 Prasanna--Venkatesh 猜想的陈述不需要导出结构，我们在此不再多说。

\subsection{导出形变环的作用}

[GV18] 构造出导出形变环之后，Galatius 与 Venkatesh 用它把 Calegari--Geraghty 方法解释成导出几何中的方法。尖点上同调的表示 $\pi$、局部化上同调 $H^*(S_{K_0},\Lambda_V)_{\mathfrak m}$、伽罗瓦参数 $\rho_\pi$ 及其模 $p$ 约化 $\overline{\rho}_\pi$、以及对偶伴随 Selmer 群 $H^1_f(\Q,M_\pi(1))$ 都如 §7.4 与 §8.1 所述。设 $S$ 是 $\pi$ 与 $K_0$ 分歧的位置集合，并假设 $p\notin S$。于是可以定义晶体导出形变环 $\mathcal R^{\overline{\rho},\mathrm{crys}}_{\OO_F[1/S]}$，它在 $S$ 处允许无限制分歧。\footnote{[Ven19] 在上一节中已承认的假设在 [GV18] 中仍然有效；其中包括 $S$ 中诸素处的局部无限制形变空间性质良好这一假设。}

Galatius--Venkatesh 对 Calegari--Geraghty 方法的解释，是用 $S_{K_0}$ 的同调而非上同调来表述的。为避免同调与上同调之间对偶的麻烦，我们假设 $V$ 是平凡表示，于是 $\Lambda_V$ 就是常局部系统 $\OO$；[GV18] 不处理更一般的情形。分次同调群 $H_*(S_{K_0},\OO)_{\mathfrak m}$ 就是 §5.5 记号下的 $H_*(\mathcal M_0)[q_0]$。下述定理是 [GV18] 的主要结果，其证明此前已概述过：

\begin{thm}{8.2.1}（[GV18, Theorem 14.1]）\label{thm:a6-05-12}
在以上各节的假设下，局部化同调 $H_*(S_{K_0},\OO)_{\mathfrak m}$ 具有分次环 $\pi_*(\mathcal R^{\overline{\rho},\mathrm{crys}}_{\OO_F[1/S]})$ 上的自由分次模结构。
\end{thm}

\ysrc{526}\subsubsection{导出 Hecke 作用与导出形变作用的对偶性}

[GV18] 一文没有给出动机猜想的新表述，这与专门讨论导出 Hecke 作用的 [Ven19] 不同。不过，当两种作用都能定义时，它们之间有典范的关联。

要点是确定 $\pi_1(\mathcal R^{\overline{\rho},\mathrm{crys}}_{\OO_F[1/S]})$ [GV18, (15.7)]：在我们的记号下，有

\begin{lem}{8.2.2}（[GV18, Lemma 15.3]）\label{thm:a6-05-13}
存在同构
\[
\pi_1\bigl(\mathcal R^{\overline{\rho},\mathrm{crys}}_{\OO_F[1/S]}\bigr)\isom\varpi_{\pi,p},
\]
它诱导出分次代数的同构
\[
\pi_*\bigl(\mathcal R^{\overline{\rho},\mathrm{crys}}_{\OO_F[1/S]}\bigr)\isom\wedge^*\varpi_{\pi,p}.
\]
\end{lem}

两种作用之间的关系不可能完全直截了当，因为导出形变环作用在同调上，而（严格整体）导出 Hecke 代数作用在上同调上。在同调与上同调的对偶之下，定理 8.1.1 中定义的 $\wedge^\bullet\varpi^*_{\pi,p}$ 在 $H^*(S_{K_0},\OO)_{\mathfrak m}$ 上的作用，变成 $H_*(S_{K_0},\OO)_{\mathfrak m}$ 上的伴随作用；$\varpi^*_{\pi,p}$ 实现为次数 $-1$ 的算子。两种作用之间的相容性取下述形式：

\begin{thm}{8.2.3}（[GV18, Theorem 15.2]）\label{thm:a6-05-14}
$\wedge^\bullet\varpi^*_{\pi,p}$ 的 DHA 作用与 $\pi_*(\mathcal R^{\overline{\rho},\mathrm{crys}}_{\OO_F[1/S]})$ 的导出形变环作用按下述意义相容。在引理 8.2.2 的同构
\[
\pi_*\bigl(\mathcal R^{\overline{\rho},\mathrm{crys}}_{\OO_F[1/S]}\bigr)\isom\wedge^*\varpi_{\pi,p}
\]
之下，设 $v\in\varpi_{\pi,p}\in\pi_1(\mathcal R^{\overline{\rho},\mathrm{crys}}_{\OO_F[1/S]})$，$v^*\in\varpi^*_{\pi,p}$，$m\in H_*(S_{K_0},\OO)_{\mathfrak m}$。则\tn{底本此处的第二个关系符号作 $\in$（$v\in\varpi_{\pi,p}\in\pi_1(\cdots)$），按上下文应为 $\subset$。（原书 p.526）}
\[
v^*\cdot v\cdot m+v\cdot v^*\cdot m=\langle v,v^*\rangle m,
\]
其中 $\langle\,,\rangle$ 是 $\varpi_{\pi,p}$ 与 $\varpi^*_{\pi,p}$ 之间自然的 $\OO$-值配对。
\end{thm}

\subsection{凝聚上同调}

\subsubsection{导出 Hecke 作用}

本节对模曲线使用经典记号。设 $d>1$ 是整数，令 $X_0(d)$ 与 $X_1(d)$ 为熟悉的模曲线，它们定义在 $\Z[\frac1d]$ 上，水平结构分别为 $\Gamma_0(d)$ 型与 $\Gamma_1(d)$ 型。我们把 $X_?(d)$（$?=0,1$）看作商
\[
\Gamma_?(d)\backslash\mathfrak H;\qquad\mathfrak H=\mathrm{GL}_2(\R)^0/\mathrm{SO}_2\cdot\R^\times,
\]
其中 $\mathrm{GL}_2(\R)^0$ 是 $\mathrm{GL}_2(\R)$ 的单位连通分支。$X_?(d)$（$?=0,1$）上以权 1 模形式为截面的线丛记作 $\omega$；$\mathrm{SO}_2$ 在 $\omega$ 的纤维（在其不动点 $h\in\mathfrak H$ 处的纤维）上的迷向表示记作 $\omega_h$。

丛 $\omega$ 是唯一一个在多于一个次数上（按 (7.3.3) 的意义）具有非平凡尖点凝聚上同调的自守线丛——换句话说，它只在次数 0 与 1 上有非平凡尖点上同调。此外，$H^0_0(X_?(d),\omega)$ 与 $H^1_0(X_?(d),\omega)$ 作为经典 Hecke 代数上的模同构，该代数记作 $\mathbb T_{?d}$；它由所有不整除 $d$ 的素数处的标准 Hecke 算子生成。对附着于所有同余子群的模曲线，这一点依然成立：单位连通分支 $\mathrm{GL}_2(\R)^0$ 有两个表示（Harish-Chandra 模），记作 $\pi_0$ 与 $\pi_1$，使得
\[
\dim H^i(\mathfrak P,\mathfrak k;\pi_i\otimes\omega_h)=1,\qquad i=0,1
\]
（记号同 (7.3.2)。）
\ysrc{527}表示 $\pi_0$ 与 $\pi_1$ 分别是全纯与反全纯的离散系列极限，而直和 $I(0):=\pi_0\oplus\pi_1$ 具有不可约 $(\mathrm{Lie}(\mathrm{GL}_2),\mathrm O_2(\R)\cdot\R^\times)$-模的结构，它限制到 $(\mathrm{Lie}(\mathrm{GL}_2),\mathrm{SO}_2\cdot\R^\times)$ 上时是 $\mathrm{GL}_2(\R)^0$ 的两个 Harish-Chandra 模之和。记号 $I(0)$ 意在提示（并且准确地提示）一点：在相差行列式绝对值的某个幂次这一无害扭之后，它可由 Borel 子群的酉特征标作抛物诱导得到。

对任意素数 $q$，令 $G_q=\mathrm{GL}(2,\Q_q)$，$K_q=\mathrm{GL}(2,\Z_q)$。设 $q$ 与 $d$ 互素，且对某个 $r>0$ 有 $q\equiv1\pmod{p^r}$。令 $X_?(d)_r$ 表示 $X_?(d)$ 到 $\Spec(\Z/p^r\Z)$ 的基变换，并类似地定义 $\omega_r$。[HV19] 定义了导出 Hecke 代数 $\mathcal H(G_q,K_q;\Z/p^r\Z)$ 在 $H^*(X_?(d)_r,\omega_r)$ 上的一个作用。该作用可以按 §7.4 的形式方式定义，但也有基于循环覆盖的具体定义。令 $X_{0?}(qd)$ 表示对应于同余子群 $\Gamma_?(d)\cap\Gamma_0(q)$ 的模曲线，并类似地定义 $X_{1?}(qd)$。只要 $q\ge5$，自然映射 $X_{1?}(qd)\to X_{0?}(qd)$ 就通过一个 \'etale 覆盖 $X_\Delta(qd)/X_{0?}(qd)$ 分解，后者是次数为 $p^r$ 的循环覆盖。这个覆盖定义一个类
\begin{equation}\tag{8.3.1}\label{eq:a6-05-9}
\mathfrak S\in H^1(X_{0?}(qd)_r;C_r),
\end{equation}
称为 Shimura 类。这里 $C_r$ 是 $(\Z/q\Z)^\times$ 的阶为 $p^r$ 的循环商。它可以与 $\Z/p^r\Z$ 等同，但并非典范地等同；我们选取一个同构（离散对数）
\begin{equation}\tag{8.3.2}\label{eq:a6-05-10}
\log:C_r\isom\Z/p^r\Z.
\end{equation}

另一方面，有两个映射 $p_i:X_{0?}(qd)\to X_?(d)$。\tn{底本此句用 $p_i$，而同段下文与命题 8.3.1 均用 $\pi_i$，疑为 $\pi_i$ 之误。（原书 p.527）}若把 $X_{0?}(qd)$ 看作三元组 $(E,C,\phi)$ 的模空间，其中 $E$ 是（广义）椭圆曲线，$C$ 是 $q$ 阶循环子群，$\phi$ 是 $d$ 处的某个水平结构，则 $\pi_1$ 是忘掉 $C$ 的映射，$\pi_2$ 是把 $E$ 换成 $E/C$ 的映射。

\begin{prop}{8.3.1}\label{thm:a6-05-15}
\begin{enumerate}[(i)]
\item 复合
\[
\begin{aligned}
H^0_0(X_?(d)_r,\omega_r)&\xrightarrow{\ \pi_1^*\ }H^0_0(X_{0?}(d)_r,\omega_r)\xrightarrow{\ \cup\mathfrak S\ }H^1_0(X_{0?}(d)_r,\omega_r)\\
&\xrightarrow{\ \pi_{2,*}\ }H^1_0(X_?(d)_r,\omega_r),
\end{aligned}
\]
（其中中间的箭头借助离散对数把 $\omega_r\otimes C_r$ 与 $\omega_r$ 等同）与一个（确切的非平凡）元素 $T^1_q\in\mathcal H^1(G_q,K_q)$ 的作用吻合。
\item $\mathcal H^1(G_q,K_q;\Z/p^r\Z)$ 在 $H^*(X_?(d)_r,\omega_r)$ 上的作用，与不整除 $qd$ 的素数处经典 Hecke 算子的作用交换。
\end{enumerate}
\end{prop}

第 (i) 部分是 [HV19] 中的等式 (3-1)，第 (ii) 部分由 Mazur 的 Eisenstein 理想论文 [Maz77] 观察到（且容易证明）。

对附着于符号差为 $(1,n-1)$ 的酉群的志村簇，S. Atanasov 在哥伦比亚大学的学位论文 [Ata22] 中构造并研究了类似的算子。这类群也有成对的（非退化）离散系列极限 $\pi_0,\pi_1$，它们在同一个（极大紧子群）$K$ 的不可约表示中取值时有 $(\mathfrak P,\mathfrak k)$-上同调。

\subsubsection{猜想}

整体 Hecke 代数 $\mathbb T_{?d}$ 在权 1 凝聚上同调空间 $H^i_0(X_?(d),\omega)\otimes\overline{\Q}_p$（$i=0,1$）上的作用可对角化，且分解中出现的特征标 $\lambda:\mathbb T_{?d}\to\overline{\Q}_p$ 与 $\mathrm{GL}(2,\mathbb A)$ 的尖点自守表示 $\pi$ 一一对应，其中 $\pi$ 的阿基米德分量 $\pi_\infty$ 同构于上面引入的不可约表示 $I(0)$。如前所述，同样的特征标在 $H^0_0$ 与 $H^1_0$ 中出现，且具相同的

\ysrc{528}重数。同样方便的是，把每个贡献到权 1 凝聚上同调的 $\pi$ 与它的规范化新形式（newform）$g_\pi$ 等同起来；后者由它的经典 $q$-展开唯一确定。在猜想的陈述中我们只写 $g=g_\pi$，并始终假设 $g$ 是水平为 $d$ 的新形式。

Venkatesh 猜想的类比在这一设定下更容易陈述，因为假设中的动机 $M_\pi$ 确实存在。由 Deligne--Serre 定理 [DS74]，每个 $\pi$（或 $g_\pi$）在 Langlands 对应的意义下对应于一个 2 维伽罗瓦表示
\[
\rho_\pi:\Gal(L/\Q)\to\mathrm{GL}(2,\overline{\Q}),
\]
其中 $L/\Q$ 是某个有限扩张，$\overline{\Q}$ 是任一代数闭域。每个这样的 $\rho_\pi$ 都是奇的：对任意复共轭 $c\in\Gal(L/\Q)$，$\rho_\pi(c)$ 的行列式等于 $-1$。反过来，Artin 猜想断言：$\Q$ 的有限扩张的伽罗瓦群的每个 2 维奇表示都具有 $\rho_\pi$ 的形式；自 Khare--Wintenberger [KW09] 关于 Serre 猜想的工作以来这一点已知，不过在多数情形下它早已由 Buzzard--Dickinson--Shepherd-Barron--Taylor [BDSBT01] 建立。此外，当 $\rho_\pi$ 的像是二面体群时，这一结果本质上归功于 Hecke。

于是对每个 $\pi$ 我们都有 2 维表示 $\rho_\pi$，它可以看作秩 2 的 Artin 动机的伽罗瓦实现，系数在一个代数整数环 $\OO$ 中；该动机是 0 维动机 $R_{L/\Q}\Spec(L)$ 的直因子。动机 $M_\pi$ 则是 $R_{L/\Q}\Spec(L)$ 的直因子，其伽罗瓦实现由 3 维伽罗瓦表示
\[
\mathrm{Ad}^0\rho_\pi:\Gal(L/\Q)\to\mathrm{GL}(3,\C)
\]
给出，即 4 维表示 $\mathrm{Ad}\,\rho_\pi$ 的无迹直和项。（一般地，$\mathrm{Ad}^0\rho_\pi$ 会通过 $L$ 的某个真子域的伽罗瓦群分解，但这不重要。）[HV19, §2.7] 解释了：在猜想的表述中，动机上同调群 $H^1_{\mathrm{mot}}(\Q,M_\pi(1))$（上面记作 $\varpi_\pi$）可以换成一个 1 维的 Stark 单位空间，记作 $U_\pi$；若写 $g=g_\pi$，则记作 $U_g$。它定义为
\begin{equation}\tag{8.3.3}\label{eq:a6-05-11}
U_g=U_\pi\isom\Hom_{\OO[G_{L/\Q}]}\bigl(\mathrm{Ad}^0\rho,\;U_L\otimes_\OO\OO\bigr).
\end{equation}
与 $E=\mathrm{Frac}(\OO)$ 张量之后它是 1 维的。我们引入一个自然映射，记作
\[
\mathrm{red}_q:U_g\to\F_q^\times\otimes_\Z\Z/p^r\Z.
\]
[HV19] 中陈述的猜想把导出 Hecke 算子 $T^1_q$ 的作用与 $U_g$ 的某个生成元的离散（模 $p^r$）对数的有理倍数联系起来。由于有理乘法与模 $p^r$ 约化不易相容，需要仔细定义输入与关系。$H^1(X_{0?}(qd)_r;C_r)$ 中的 Shimura 类通过 (8.3.2) 与 Zariski 层包含 $\Z/p^r\Z\hookrightarrow\OO_r$ 的复合，定义一个非平凡类 $\mathfrak S_{\OO}\in H^1(X_{0?}(qd)_r;\OO_r)$，其中 $\OO_r$ 在这里表示 $X_{0?}(qd)_r$ 的结构层。$\mathrm{GL}_2$ 及其内型所特有的一个现象是：类 $\mathfrak S_{\OO}$ 是 $H^1(X_{0?}(qd),\OO)$ 中（某个）Eisenstein 尖点形式 $f^{\#}_E$ 的复共轭的类模 $p^r$ 约化，这里「Eisenstein 尖点形式」取 [Maz77] 的意义。\footnote{Atanasov 在 [Ata22] 中证明：对符号差为 $(1,n-1)$（$n>2$）的酉群所附着的志村簇，Shimura 类的类似物不能提升到特征 0。}至少当 Eisenstein 尖点形式空间是

\ysrc{529}一维时，这使我们能指定一个不变量，用来衡量 $T^1_q$ 的作用。定义 $g^*$ 为对应于对偶自守表示 $\pi^\vee$ 的权 1 新形式。令 $g^{*,(q)}$ 为空间 $\pi^\vee$ 中的「旧形式」，在经典记号下它就是上半平面上的解析函数 $z\mapsto g^*(qz)$。由命题 8.3.1，$T^1_q(g)\in H^1(X_{0?}(d)_r,\omega_r)$ 与 $g$ 具有相同的 Hecke 特征值。因此在 Serre 对偶配对
\[
\langle\,,\rangle:H^0_0(X_{0?}(d)_r,\omega_r)\otimes H^1_0(X_{0?}(d)_r,\omega_r)\to H^1_0(X_{0?}(d)_r,\omega_r^{\otimes2})\to\OO_r
\]
（最后箭头是迹映射）中，映射 $\bullet\mapsto\langle\bullet,T^1_q(g)\rangle$ 通过到水平 $qd$ 的 $g=g_\pi$-特征空间上的投影分解。现在令 $f^{\#}_E=a_1\cdot f_E$，其中 $f_E$ 是规范化的 Eisenstein 尖点形式（首项系数为 1；常数 $a_1$ 本质上是 [HV19] 中所谓 Merel 常数的倒数）。于是 $T^1_q(g)$ 由配对 $\langle gg^{*,(q)},f_E\rangle$ 确定，而后者可以等同于权 2 形式 $gg^{*,(q)}$ 与 $f_E$ 的 Petersson 配对；由于 $f_E$ 的水平为 $q$，$g^*$ 必须换成旧形式（该配对发生在命题 8.3.1 中右起第二个群里）。它还可以与三重积 $L$-函数 $L(s,\pi\times\pi^\vee\times\pi(f_E))$ 中心值的平方根联系起来，其中记号 $\pi(f_E)$ 本身已说明含义。

[HV19] 中表述的动机猜想归结为一个相差整数因子的等式：

\begin{conj}{8.3.2}\label{thm:a6-05-16}
存在整数 $m=m_g\ge1$ 与 $u_g\in U_g$，使得对上述所有素数 $q$ 与 $p$ 都有
\[
m\cdot\langle gg^{*,(q)},\mathfrak S\rangle=\log(\mathrm{red}_q(u_g)).
\]
\end{conj}

这里两端都按 (8.3.2) 中 $\log$ 的选取线性地变化。$m$ 与 $q$、$p$ 的无关性在这一陈述中是本质的，否则陈述将空洞无物。

\subsubsection{结果}

权 1 特征形式 $g=g_\pi$ 可以按相应的（奇）2 维伽罗瓦表示 $\rho_\pi$ 的像分类：一类像是二面体群——由二次扩张 $K/\Q$ 的一个特征标 $\psi_1$ 诱导而来；另一类是「奇异」情形，其射影化的像是 $A_4$、$S_4$ 或 $A_5$。在二面体情形，可以确定 Stark 单位群 $U_\pi$。当 $K$ 为虚域——即 $g$ 或 $\pi$ 是 CM 型——时，$U_\pi$ 由椭圆单位组成；当 $K$ 为实域——此时 $g$ 或 $\pi$ 是 RM 型——时，$U_\pi$ 由 $K$ 本身的基本单位生成。设 $D$ 是 $K$ 的判别式。利用 $U_\pi$ 的显式确定，[DHRV22] 证明了如下定理：

\begin{thm}{8.3.3}\label{thm:a6-05-17}
若 $K$ 为虚域，则假设 $D$ 是奇素数且 $\psi_1$ 非分歧。若 $K$ 为实域，则假设 $D$ 是奇数且 $\psi_1$ 的导子整除 $K/\Q$ 的差积。此时猜想 8.3.2 对 $g_\pi$ 成立。
\end{thm}

对 $D$ 与 $\psi_1$ 的导子的这些限制当然不是必需的；R. Zhang 在哥伦比亚大学的学位论文 [Zha23] 中放松了它们，其结果接近最优。定理 8.3.3 的证明方法关键地依赖于把 $g_\pi$ 实现为显式的 theta 函数。D. Marcil 在某些 $A_4$ 情形给出了该猜想引人注目的数值验证 [Mar21]。就目前情况看，在奇异情形证明无条件结果似乎不太可能，尽管可以设想该猜想会以某种方式作为 Stark 猜想 [Sta71,Sta75,Sta76,Sta80] 的推论而得到。

\ysrc{530}\subsection{Venkatesh 的猜想、$p$-进上同调与不变量 $\ell_0$}

$p$-进模形式专家早在 1990 年代就遇到了不变量 $\ell_0$。Hida 关于普通模形式（或自守上同调类）的理论，对 $\Q$ 上的约化群 $G$ 构造了一个 $p$-进解析空间，其点参数化 $p$-进模形式；该空间可看作形式谱 $\Spf(\mathbb T)$ 的一般纤维 $\mathcal X_{\mathrm{ord},G}$，其中 $\mathbb T$ 是 Hida 的「大 Hecke 代数」。环 $\mathbb T$ 是 Iwasawa 代数 $\Lambda$ 上的代数，$\Lambda$ 对应（$G$ 的 $p$-进点的极大紧子群所对应的）极大环面。在 Hida 关于模曲线的原始理论及其到志村簇的推广中，$\mathbb T$ 是有限平坦 $\Lambda$-代数，而对应于经典（复）Hecke 代数特征值的 $\mathcal X_{\mathrm{ord}}$ 的点在 $\mathcal E_{\mathrm{ord}}$ 中稠密。\tn{底本此处称这些点在 $\mathcal E_{\mathrm{ord}}$ 中稠密，但 $\mathcal E_{\mathrm{ord}}$ 未另行定义，疑为 $\mathcal X_{\mathrm{ord}}$ 之误。（原书 p.530）}一般地，Hida 在 [Hid98] 中猜想：$\mathcal X_{\mathrm{ord},G}$ 在 $\Spf(\Lambda)$ 的刚性一般纤维中的余维数恰为 $\ell_0$。

Urban 在论文 [Urb11] 中作了类似的猜想；该文中他构造了有限斜率的模形式的 $p$-进族，由特征簇 $\mathcal X_G$ 参数化，这是椭圆模形式的 Coleman--Mazur 特征曲线的多种高维推广构造之一。在有限斜率情形，特征簇纯粹是刚性解析对象，Hida 的 Hecke 代数换成整体截面环 $\OO_{\mathcal X_G}$ 的一个 Noether 子环（我们也记作 $\mathbb T$），而 $\Spf(\mathbb T)$ 到 $\Spf(\Lambda)$ 的映射换成一个态射 $\mathrm{pr}:\mathcal X_G\to\mathcal W$。这里权空间 $\mathcal W$ 仍是维数等于 $G$ 的绝对秩的刚性环面，Iwasawa 代数换成点 $w\in\mathcal W$ 处的完备（Noether）局部环 $\Lambda_w$。固定点 $x\in\mathcal X_G$，设 $\mathrm{pr}(x)=w$，则映射 $\mathrm{pr}$ 使完备化 $\mathbb T_x$ 成为有限 $\Lambda_w$-代数。

Hida 与 Urban 都把他们的余维数猜想称为 Leopoldt 猜想的非阿贝尔类比；Leopoldt 猜想说的是数域中整体单位的 $p$-进乘法独立性。数十年来，Leopoldt 猜想一直是代数数论的百慕大三角，引诱许多最杰出的专家在追求有希望但最终徒劳的证明尝试中耗费宝贵年华。换句话说，最好别去试着证明 Hida 与 Urban 的余维数猜想。不过，在 Venkatesh 的动机框架中重新解释这些余维数猜想是很自然的。我们简要讨论两种这样的解释。Hansen 与 Thorne 的文章 [HT17] 把特征簇当作伽罗瓦表示的参数空间来处理；而 Khare 与 Ronchetti 的文章 [KR23b] 提供了证据，说明 §8.4.2 中引入的导出钻石代数作用在 Hida 的普通族上时，是 Venkatesh 在（驯顺的）Taylor--Wiles 素处的导出 Hecke 代数的野 $p$-进类比。在这两种构造中，Venkatesh 动机猜想中出现的动机上同调群 $H^1_{\mathrm{mot}}(\Q,M_\pi(1))$ 都通过其 $p$-进实现、即对偶 Selmer 群出现。

\subsubsection{Hansen--Thorne 构造}

[HT17] 一文只讨论附着于群 $G=\mathrm{GL}(n)_\Q$ 的特征簇。这些特征簇按 Hansen 在 [Han17] 中的方法构造，遵循 Ash 与 Stevens 的拓扑进路。由 Harris--Lan--Taylor--Thorne [HLTT16] 的工作以及 Scholze [Sch15b] 的工作可知：对应于尖点上同调自守表示 $\pi$ 的经典点 $x\in\mathcal X:=\mathcal X_{\mathrm{GL}(n)}$ 配有一个半单 $p$-进伽罗瓦表示
\begin{equation}\tag{8.4.1}\label{eq:a6-05-12}
\rho_\pi:\Gal(\overline{\Q}/\Q)\to\mathrm{GL}(n,L),
\end{equation}
